\section{The full quadratic germ and its physical return}
\label{sec:normalization}

We first separate the affine and positive-time changes from the five
essential coefficients. We then construct the return through both
singular points, including the two regular closing maps. The resulting
displacement is written in an affine section height, which will be used
throughout the zero count.

\subsection{Reduction to five parameters}

The following five-parameter family is transverse to the seven affine
and positive-time directions:
\begin{equation}\label{eq:slice}
 H_\theta:\quad
 \begin{cases}
 \dot X=AX-Y+X^2+(\mu_2+\mu_3)XY+\mu_1Y^2,\\
 \dot Y=CX+X^2+XY+\mu_3Y^2,
 \end{cases}\qquad
 \theta=(A,C,\mu_1,\mu_2,\mu_3)
\end{equation}
near \(\theta_0=(1,1,0,0,0)\).

\begin{lemma}\label{lem:slice}
There are neighborhoods \(\mathcal U\subset\mathcal Q_2\) of \(H_0\) and
\(\mathcal W\) of \((I,0,1,\theta_0)\) such that every
\(H\in\mathcal U\) has a unique representation with parameters in
\(\mathcal W\),
\begin{equation}\label{eq:affine-action}
 H(z)=a_{\rm t}L^{-1}H_\theta(Lz+c),
 \qquad (L,c,a_{\rm t},\theta)\ \hbox{near }(I,0,1,\theta_0),
\end{equation}
with \(a_{\rm t}>0\).  The resulting maps
\(H\mapsto \mathfrak A(H):=(L,c,a_{\rm t})\) and
\(H\mapsto\theta(H)\) are
analytic on \(\mathcal U\), and satisfy
\(\mathfrak A(H_0)=(I,0,1)\) and \(\theta(H_0)=\theta_0\).
\end{lemma}
\begin{proof}
The derivative of the coefficient map has columns
\[
 DH_0h-Dh\,H_0,\quad
 h=(1,0),(0,1),(X,0),(Y,0),(0,X),(0,Y),
\]
followed by \(H_0\) and the five vectors
\[
 (X,0),\quad(0,X),\quad(Y^2,0),\quad(XY,0),\quad(XY,Y^2).
\]
In monomial order \((1,X,Y,X^2,XY,Y^2)\) in each component,
the determinant is \(3\). The analytic inverse function theorem
proves the claim; positivity of \(a_{\rm t}\) persists near \(1\).
\end{proof}

For \(\mathfrak A=(L,c,a_{\rm t})\), denote by
\(T_{\mathfrak A}\) the extension of
\(z\mapsto Lz+c\) to the closed hemisphere, obtained by normalizing
\((X,Y,W)\mapsto(L(X,Y)^{\mathsf T}+cW,W)\).  After shrinking
\(\mathcal U\) to have compact closure in the inverse-function
neighborhood, these extensions and their inverses have a common Lipschitz
constant \(1\le L_*<\infty\).  Moreover,
\begin{equation}\label{eq:slice-graphic-motion}
 d_{\rm H}\bigl(T_{\mathfrak A(H)}\Gamma,\Gamma\bigr)\longrightarrow0
 \qquad\text{as }H\longrightarrow H_0,
\end{equation}
because \(\mathfrak A(H)\to(I,0,1)\).

\subsection{The two singular points}

In the chart at the positive \(Y\)-direction, use
\(v=X/Y,\ z=1/Y\). Multiplying the vector field by \(z>0\) preserves
the direction of its orbits in the finite plane and gives
\begin{equation}\label{eq:raw}
 \dot v=D_\theta(v,z)=P_\theta(v)+Q_\theta(v)z,\qquad
 \dot z=-zG_\theta(v,z)=-z\{R_\theta(v)+S_\theta(v)z\},
\end{equation}
where
\begin{equation}\label{eq:polynomials}
 \begin{aligned}
 P(v)&=\mu_1+\mu_2v-v^3,& Q(v)&=-1+Av-Cv^2,\\
 R(v)&=\mu_3+v+v^2,& S(v)&=Cv .
 \end{aligned}
\end{equation}
The physical side is \(z>0\); the equator \(z=0\) is invariant.
We always use the positive-time chart field on this side.
At the organizing parameter \(\theta_0=(1,1,0,0,0)\), the chart is
\begin{equation}\label{eq:infinity-chart-field}
 \dot v=-v^3-(1-v+v^2)z,\qquad
 \dot z=-z(v+v^2+vz).
\end{equation}
The linearization is \(\bigl(\begin{smallmatrix}0&-1\\0&0\end{smallmatrix}\bigr)\).
Eliminating the analytic unit \(1-v+v^2\) gives the nullcline
\(z=-v^3/(1-v+v^2)\), on which
\begin{equation}\label{eq:infinity-multiplicity-four}
 \left.\dot z\right|_{\dot v=0}
 =\frac{v^4}{(1-v+v^2)^2}.
\end{equation}
Here \((1+v)(1-v+v^2)=1+v^3\); the reduced zero has order four,
the multiplicity of the nilpotent singularity at infinity.
The organizing arc has equation
\begin{equation}\label{eq:organizing-arc}
 z_\Gamma(v)=\sqrt{1+v^2}-1
\end{equation}
and is traversed with decreasing \(v\); its expansion
\(z_\Gamma(v)=v^2/2+O(v^4)\) gives the common equatorial tangent.
Choose a small fixed \(\delta>0\) and the sections
\(\Sigma_\pm=\{v=\pm\delta\}\), with small height intervals around this arc.
Take the height intervals small enough that \(-D\) has a positive lower
bound on their closures.
On the fixed small chart, a coefficient shrink gives
\begin{equation}\label{loc:signs}
 Q<-q_0<0,\qquad G_v=1+2v+Cz>g_0>0,\qquad
 P_\theta(\delta)<0,\quad G_\theta(-\delta,z)<0.
\end{equation}

The small zero of the transverse eigenvalue and the first two
equatorial coefficients there are
\begin{equation}\label{eq:center-parameters}
 \gamma=\frac{-1+\sqrt{1-4\mu_3}}2,\qquad
 \alpha_0=P(\gamma),\qquad \alpha_1=P'(\gamma).
\end{equation}
The Jacobian with respect to \((\mu_3,\mu_1,\mu_2)\) at \(\theta_0\)
is \(-1\), so these, together with \(A,C\), are parameter coordinates.
The translated polynomials and their rescalings are introduced in
Section~\ref{sec:coefficient-regimes}. The five raw coefficients
\(A-1,C-1,\gamma,\alpha_0,\alpha_1\) vary independently.

For the finite saddle-node, set \(u=X+1,\ W=Y+u\).
At \(\theta_0\),
\[
 \dot u=u^2-W,\qquad \dot W=u^2+(u-2)W .
\]
The zero set of the second component is an analytic graph
\(W=\psi(u,\theta)\), used here only to locate equilibria.
The first component on this graph has a Weierstrass factor
\[
 e(u,\theta)(u^2+b_1(\theta)u+b_0(\theta)),\qquad
 e(0,\theta_0)=1/2 .
\]
Indeed, at \(\theta_0\) it equals \(u^2(1-u)/(2-u)\).
Differentiation gives
\[
 b_0=C-A+O(|\theta-\theta_0|^2),\qquad
 b_1=A-1+\mu_2+\mu_3+O(|\theta-\theta_0|^2).
\]
Consequently \(K=b_0-b_1^2/4\) is a discriminant coordinate with
\(dK=dC-dA\) at the organizing field.
We also write \(K_{\rm fin}=K\) for this finite-passage discriminant.

The organizing analytic germ has eigenvalues $0,-2$ and a nonzero
quadratic weak coefficient $1/2$. Thus its multiplicity is two, the
case $\mu=1$ of the saddle--node normal form. Fix a finite joint
differentiability order \(N\ge12\) before restricting the coefficients.
The ordinary section maps and their inverses will be used through phase
order three, while division by the discriminant costs at most one
parameter derivative. The intrinsic root-collar kernels use phase order
five, axis order six and selected mixed order seven; these estimates and
the strong projection estimates are obtained from the analytic chart
field. Choose the order for the normal-form theorem now, increasing it
if necessary before choosing any neighborhood. This finite surplus
accommodates coefficient differentiation and Hadamard division, so a
jointly smooth normalizer of all orders is unnecessary.
The finite-smooth multiparameter unfolding normal form
\cite[Theorem~3 and formula~(2.2), October 2000 author copy p.~6]{DIR02} gives,
by a joint \(C^N\) orbital equivalence preserving positive time,
\begin{equation}\label{eq:finite-normalform}
 \xi'=\eps_f+\frac{\xi^2}{1+b\xi},\qquad \eta'=-\eta .
\end{equation}
Here the parameter in the cited theorem is the complete vector
\(\theta-\theta_0\). For \(\mu=1\), its polynomial \(P(\theta,\xi)\)
is the constant \(\eps_f\). The reduction divides by the nonzero strong
rate, which is positive at the organizing parameter and remains so
nearby; it thus gives the stated positive-time equivalence. The theorem
imposes no rank condition on the unfolding coefficients.
The linear calibration
\(t_{\rm nf}=2t,\ \xi=(X+1-Y)/8+O(|(X+1,Y)|^2)\) gives
\(\eps_f=(C-A)/16+O(|\theta-\theta_0|^2)\).
The double-equilibrium loci in the two descriptions coincide.
In parameter coordinates \((K,\vartheta)\), Hadamard division gives
\[
 \eps_f(K,\vartheta)=K h(K,\vartheta),\qquad
 h(K,\vartheta)=\int_0^1\partial_K\eps_f(tK,\vartheta)\,dt,
 \qquad h(\theta_0)=1/16.
\]
Continuity makes \(h\) a positive unit on a smaller coefficient
neighborhood.

For fixed internal sections \(\xi=\pm\rho\), with \(1+b\xi>1/2\),
the passage exists when \(K>0\) and is exactly
\begin{equation}\label{eq:finite-multiplier}
 \eta_{\rm out}=m_{\rm fin}\eta_{\rm in},\quad
 m_{\rm fin}=e^{-T_f},\quad
 T_f=\int_{-\rho}^{\rho}
       \frac{d\xi}{\eps_f+\xi^2/(1+b\xi)}
     =\frac{\pi}{\sqrt{\eps_f}}+O(1).
\end{equation}
We write \(m=m_{\rm fin}\) in the completed return below.
Completing the square in \(\xi^2+b\eps_f\xi+\eps_f\) proves the last
estimate. Differentiating its arctangent and logarithmic expressions
also shows that each fixed finite parameter derivative of \(m\) is
bounded by \(CK^{-M}e^{-c/\sqrt K}\). Thus \(m\) extends flatly
by zero to \(K=0\), to the chosen differentiability order.
If \(K\le0\), an orbit entering the selected channel approaches the
saddle-node or the left equilibrium and cannot exit through
\(\xi=\rho\). The strong sides of this box point inward.

\subsection{The exact displacement}

Let \(T_\theta:\Sigma_+\to\Sigma_-\) be the infinity transition when
it exists in positive time. The rest of the return consists of two
regular transitions and the finite passage
\eqref{eq:finite-multiplier}; see Figure~\ref{fig:graphic-sections}.

\begin{lemma}\label{lem:closing}
There are uniformly regular increasing maps \(C_\theta,E_\theta\)
on fixed section intervals such that the finite-side return is
\(q\mapsto C_\theta(mE_\theta(q))\). Their derivatives and inverse
derivatives are bounded through the chosen finite order on buffered
section intervals.
If \(a_\theta=C_\theta(0)\), \(f_\theta=C_\theta-a_\theta\),
and \(d_\theta=E_\theta\), then in the affine coordinate
\(x=z_{\rm in}-a_\theta\) the cycle equation is
\begin{equation}\label{eq:return}
 \Delta_\theta(x)=x-m\Pi_\theta(x)=0,\qquad
 \Pi_\theta(x)=\mathcal C_{\theta,m}(T_\theta(a_\theta+x)),
\end{equation}
where
\begin{equation}\label{eq:divided-closing}
 \mathcal C_{\theta,m}(q)=d_\theta(q)
       \int_0^1 f_\theta'(tm d_\theta(q))\,dt .
\end{equation}
In the output coordinate \(t\), the equivalent equation is
\begin{equation}\label{eq:inverse-return}
 S_\theta(t)-a_\theta-f_\theta(mt)=0,\qquad
 S_\theta=(d_\theta\circ T_\theta)^{-1}.
\end{equation}
\end{lemma}
\begin{proof}
The parameter-dependent flow-box theorem on the two compact regular
segments gives positively oriented \(E_\theta,C_\theta\), with bounded
derivatives and reciprocal slopes on smaller fixed intervals.
The integral in \eqref{eq:divided-closing} extends
\(f_\theta(md_\theta(q))/m\) to \(m=0\), with positive slope
\(f_\theta'(md_\theta)d_\theta'\). Substitution proves both equations
and preserves isolated zeros and their multiplicities.
\end{proof}

At a zero, boundedness of \(\mathcal C_{\theta,m}\) gives
\(|x|\le Cm\).
In the inverse equation, the nonlinear closing term has an additional
factor \(m^2\):
\begin{equation}\label{eq:closing-quadratic}
 f_\theta(mt)=mf_\theta'(0)t+
 m^2t^2\int_0^1(1-u)f_\theta''(umt)\,du .
\end{equation}
The affine input keeps this small quadratic closing term explicit in
the later derivative comparisons.

The separate value-and-slope normalization
\eqref{eq:reference-normalization} is applied to the actual input height
and completed output height. It retains their two centers and gains;
only the input normalization is affine in the physical incoming height.

\subsection{Order and uniform transverse sections}
\label{sec:localization}

We first bound the through trajectories from above by an invariant graph.
Choose \(z_{\rm r}=\kappa\delta\), with
\(\delta\ll\kappa\ll1\). On
\(z_{\rm r}/2\le z\le2z_{\rm r}\), a coefficient shrink gives
\(|P|\le z_{\rm r}/8,\ Q\le-1/2,\ |G|\le C\delta\).
The orbit through \((\delta,z_{\rm r})\) consequently has
\(D\le-z_{\rm r}/8\) and \(|dz/dv|\le C\delta\).
It crosses the chart inside this band, by a first-exit argument.
Its graph \(\overline z_\theta(v)\) is the upper boundary; all
selected section windows are strictly below it.

\begin{lemma}\label{lem:loc-through}
The heights on \(\Sigma_+\) which reach \(\Sigma_-\) in the chart
below this upper boundary form an interval
\begin{equation}\label{loc:through-interval}
 \mathcal T_\theta=(\Theta_\theta,\overline z_\theta(\delta)] .
\end{equation}
Every such orbit has \(D<0\). Both \(T_\theta\) and
\begin{equation}\label{loc:speed-margin}
 \sigma_\theta(h)=\min_{|v|\le\delta}
          \{-D_\theta(v,z_\theta(v;h))\}
\end{equation}
are strictly increasing functions of the input \(h\).
\end{lemma}
\begin{proof}
Every physical input has \(D<0\) on the right section by
\eqref{loc:signs}. On the nullclines,
\begin{equation}\label{loc:nullcline}
 \dot D|_{D=0}=-zQG,\qquad \dot G|_{G=0}=G_vD .
\end{equation}
A first crossing from \(D<0\) to \(D>0\) requires \(G>0\);
equality would be an equilibrium. The quadrant \(D>0,G>0\)
is forward invariant by these identities, so this orbit cannot
subsequently reach the left section. A first terminal nullcline hit on
\(v=-\delta\) is also impossible: there \(\dot D<0\) by
\eqref{loc:signs}. A through orbit therefore has \(D<0\) everywhere
and solves
\begin{equation}\label{loc:graph-equation}
 z_v=-zG/D .
\end{equation}
Scalar uniqueness strictly orders these graphs. A graph above a through
graph has smaller \(D\), since \(Q<0\), and cannot cross either
\(D=0\) or the invariant upper boundary. It also passes through;
openness proves the interval assertion and ordering proves that \(T\)
increases. The positive difference of the two \(-D\)-values has a
positive minimum on the compact \(v\)-interval, so \(\sigma\) increases.
\end{proof}

The lower endpoint is a stopped limiting passage. A \emph{gap} is the
positive height difference from it on a specified transverse cut,
distinct from the central height.

\begin{lemma}\label{lem:loc-stopped-passage}
There are \(0<r<\delta\) and \(d_0,c_r>0\), independent of the
through orbit and of the sufficiently small coefficient, such that
every through orbit with endpoints in the selected outer windows and
\(\sigma_\theta<d_0\) crosses \(v=r\) and \(v=-r\) once, in order.
On \(r\le |v|\le\delta\) it has \(-D\ge c_r\).
The two outer transitions consequently have uniformly bounded
derivatives and inverse derivatives of any prescribed finite order
on slightly larger fixed domains.
\end{lemma}
\begin{proof}
The incidence set
\(\mathcal E_\theta=\{D=0,\ zG=0\}\) at \(\theta_0\) meets the
small closed chart only at the origin. Compactness puts the whole
nearby incidence set in \(|v|<r/3\).
Choose a neighborhood \(\mathcal O\) of it in \(|v|<2r/3\).
The endpoint lower bounds persist on the selected outer windows.
At an interior minimum of \(-D\) on a through graph,
\begin{equation}\label{loc:minimizer-identity}
 D\,D_v-QzG=0 .
\end{equation}
Transversality excludes a minimum near the fixed outer sections.
As \(\sigma_\theta\to0\), \eqref{loc:minimizer-identity} gives
\(|zG|\le C\sigma_\theta\); hence the minimizing point must enter
\(\mathcal O\). This proves the first assertion for some \(d_0\).
If the outer speed had no positive lower bound, a limiting orbit
would have a point with \(D=0\), \(D\le0\) on both time sides,
and thus \(\dot D=0\). By \eqref{loc:nullcline} this point lies
in \(\mathcal E\), a contradiction. The decreasing coordinate \(v\)
now gives all ordered hits. The hitting-time theorem on these
uniformly transverse compact slabs proves the derivative bounds.
\end{proof}

Fix nested section windows, whose interiors contain the corresponding
intersections of the organizing graphic, and speed bounds
\begin{equation}\label{loc:nested-buffers}
 K_\pm\Subset K_\pm^\sharp\Subset\Sigma_\pm,\qquad
 0<d_{\rm in}<d_{\rm out}<d_0 .
\end{equation}
The inner and outer sets specified by
\[
 h\in K_+,\quad T_\theta(h)\in K_-,\quad \sigma_\theta(h)\le d_{\rm in},
\]
and respectively by the interiors of \(K_\pm^\sharp\) and
\(\sigma_\theta(h)<d_{\rm out}\), are intervals, denoted by
\(J_\theta^{\rm in}\subset J_\theta^{\rm out}\).
All intervening transverse sections have nested windows as well.
The central height
\begin{equation}\label{loc:central-height}
 h_\theta(h)=z_\theta(\gamma(\theta);h)
\end{equation}
is positive and strictly increasing on \(J_\theta^{\rm out}\).
For any fixed raw height cap \(h_*>0\), define the affine physical
receiver
\begin{equation}\label{loc:analytic-receiver}
 I_\theta(h_*)=
 \{x:\ a_\theta+x\in J_\theta^{\rm out},
       \ 0<h_\theta(a_\theta+x)<h_*\}.
\end{equation}
It is an interval, possibly empty. Thus the whole interval between any
finite collection of its inputs satisfies the same input, inverse-output
and speed restrictions. These geometric restrictions define the domain
before any zero estimate is applied, and keep it on the through side of
the stopped boundary.

The estimates in the following sections choose a positive cap
\(h_{\rm out}\) in \eqref{loc:analytic-receiver} with strict buffers.
Fix this cap and the nested windows before the final independent raw
coefficient restrictions and the multiplier bound. Subsequent restrictions
retain these caps and section bounds. After these choices, the following
capture lemma places all nearby complete cycles below a smaller cap
\(0<h_{\rm in}<h_{\rm out}\).

\begin{lemma}\label{lem:loc-capture}
For every fixed set of these admissible nested windows and speed bounds in
\eqref{loc:nested-buffers}, and every \(h_{\rm in}>0\), there
are a coefficient neighborhood \(\mathcal V_0\) of \(\theta_0\) and
\(\eps_0>0\) such that, for every \(\theta\in\mathcal V_0\), every periodic
orbit \(C\) of \(H_\theta\) with
\(d_{\rm H}(\iota_0(C),\Gamma)<\eps_0\) follows the selected
positive-time circuit and has \(K>0\).
It meets the return section once in its prime period.
Its infinity arc has input in \(J_\theta^{\rm in}\), all selected
transverse hits in the inner windows, and
\(0<z(\gamma)<h_{\rm in}\).
\end{lemma}
\begin{proof}
Surround the two singularities by small blocks and the two compact
regular branches by flow boxes, with nested transverse windows.
In the finite normal-form box take weak faces \(\xi=-\rho\) and
\(\xi=\rho\), with transverse windows about the two organizing hits.
The weak speed at these fixed faces is positive after a coefficient
shrink. Thus they are respectively entrance and exit faces. The strong
faces \(\eta=\pm\eta_0\) point inward, since \(\eta'=-\eta\),
and are chosen away from the organizing arc. Remove from the two weak
faces the parts outside the transverse windows. These parts and the
strong faces are compact and disjoint from \(\Gamma\).

The infinity block has transverse faces \(v=\delta\) and
\(v=-\delta\), the invariant equator below, and the invariant graph
\(\overline z_\theta\) above. At the two selected face windows the
speed \(D\) is uniformly negative. Its upper face and the portions
of the transverse faces outside these windows stay a positive distance
from \(\Gamma\). The two regular flow boxes join the finite exit to
the right infinity entrance and the left infinity exit to the finite
entrance, in positive time. Their lateral faces have the same positive
distance property. Choose the Hausdorff radius below all these boundary
distances, which persist for nearby coefficients. A close cycle must
visit both regular branch boxes, because it approximates the complete
graphic. It can leave each visited box only through the selected next
face; the equator cannot be crossed by uniqueness. These entrances and
exits therefore force the finite channel, one regular branch,
\(\Sigma_+\to\Sigma_-\), and the other regular branch in this order.
The finite channel has no exit for \(K\le0\), so \(K>0\);
Lemma~\ref{lem:loc-through} gives \(D<0\) on the infinity passage.

Every intersection of the cycle with the selected connected segment
of \(\Sigma_+\) lies in its inner window and has the same transverse
orientation. A prime periodic orbit of a planar vector field is a
Jordan curve. Along a connected transverse segment its intersections
alternate between entering and leaving its bounded complementary
component. Hence these equally oriented intersections consist of one
point. The same circuit then contains one infinity arc, and its height
is one isolated zero of \eqref{eq:return} when the cycle is isolated.

If the uniform inner bounds failed, take
\(\theta_n\to\theta_0\) and \(C_n\) with
\(d_{\rm H}(\iota_0(C_n),\Gamma)\to0\).
Let \(A_n\) be the unique infinity-block subarc selected above and \(A_0\)
the complete organizing arc in that block.
Every Hausdorff limit of \(A_n\) is contained in \(A_0\).
Conversely every interior point of \(A_0\) is approached by
\(\iota_0(C_n)\), and the selected circuit forces the approaching
points to be on \(A_n\); closure includes the endpoints.
Thus \(d_{\rm H}(A_n,A_0)\to0\).
For a common modulus of continuity \(\omega\),
\[
 \left|\min_{A_n}(-D_{\theta_n})-\min_{A_0}(-D_{\theta_0})\right|
 \le \|D_{\theta_n}-D_{\theta_0}\|_\infty
          +\omega(d_{\rm H}(A_n,A_0)).
\]
Since the minimum on \(A_0\) is zero, \(\sigma_{\theta_n}(A_n)\to0\).
The hits of \(v=\gamma(\theta_n)\) converge to \((0,0)\), the unique
hit of \(v=0\) on \(A_0\), and all transverse hits enter the inner
windows by continuity. This contradiction supplies one coefficient
neighborhood and one Hausdorff radius.
\end{proof}

\subsection{Transferring a slice bound}

The same geometric transfer will be used for the full coefficient cover
and for its ratio-restricted parts.

\begin{lemma}\label{lem:slice-bound-transfer}
Let $\mathcal U_5$ be a neighborhood of $\theta_0$, let
$\mathcal B$ be any subset of the slice parameter space, and fix
$\eps_5>0$ and an integer $N_{\rm cyc}\ge0$.
Suppose every $\theta\in\mathcal U_5\cap\mathcal B$ has at most
$N_{\rm cyc}$ distinct isolated cycles at whole-graphic Hausdorff
distance less than $\eps_5$ from $\Gamma$.
There are one open neighborhood $\mathcal V$ of $F_*$ in
$\mathcal Q_2$ and one $\eps>0$ such that every $F\in\mathcal V$
with $\theta(\mathscr C_0(F))\in\mathcal B$ has at most
$N_{\rm cyc}$ distinct isolated periodic orbits $C$ satisfying
$d_{\rm H}(\iota_*(C),\Gamma_{I4})<\eps$.
The transport preserves the complete
marking, positive time and return-zero multiplicities. The maps and
their inverses also preserve whole-orbit Hausdorff convergence to the
marked graphics as the fields tend to $F_*$.
\end{lemma}
\begin{proof}
For $H=\mathscr C_0(F)$ use the analytic decomposition in
Lemma~\ref{lem:slice} for this field. Enlarge $L_0$ to bound both $S$ and
$S^{-1}$, and use the common bound $L_*$ for
$T_{\mathfrak A(H)}$ and its inverse. Continuity gives an open
neighborhood $\widehat{\mathcal V}$ of $H_0$ on which
\begin{equation}\label{eq:lift-neighborhood}
 \theta(H)\in\mathcal U_5,\qquad
 \eta(H):=d_{\rm H}(T_{\mathfrak A(H)}\Gamma,\Gamma)<\eps_5/2.
\end{equation}
Set
\begin{equation}\label{eq:original-coefficient-neighborhood}
 \mathcal V=\mathscr C_0^{-1}(\widehat{\mathcal V}),
\end{equation}
an ordinary open neighborhood in all twelve coefficients, and take
\begin{equation}\label{eq:original-collar-radius}
 \eps=\frac{\eps_5}{2L_*L_0}.
\end{equation}
The fixed change and $w=LZ+c$ send a cycle $C$ of this same field
to a cycle $\widetilde C$ of $H_{\theta(H)}$, with
\[
 \iota_0(\widetilde C)=T_{\mathfrak A(H)}S(\iota_*(C)).
\]
Writing $d_{\rm original}=d_{\rm H}(\iota_*(C),\Gamma_{I4})$
and $d_{\rm slice}=d_{\rm H}(\iota_0(\widetilde C),\Gamma)$,
we have
\begin{align}
 d_{\rm slice}&\le L_*L_0d_{\rm original}+\eta(H)<\eps_5,
                                      \label{eq:two-map-hausdorff}\\
 d_{\rm original}&\le L_0L_*(d_{\rm slice}+\eta(H)).
                              \label{eq:two-map-hausdorff-reverse}
\end{align}
The affine diffeomorphisms and positive time factors biject cycles,
preserving distinctness, isolation, the two branches and their marked
order. The induced transverse return maps are locally conjugate, so
zero multiplicities are preserved as well. The first inequality applies
the slice bound; both inequalities and $\eta(H)\to0$ give the
convergence assertion as the fields tend to the organizing field.
\end{proof}
